%% Chapter 8 -----------------------------------------------------------------
\chapter{Type II Burst Analysis}
\label{ch:type-ii}
\index{Type II burst}

Type~II solar radio bursts are slowly drifting emission lanes produced by shock waves
traveling outward through the corona. The plasma frequency falls with height, so the
drift of the emission to lower frequencies traces the shock's motion. The Type~II tools
fit the burst backbone and convert its frequency and drift into a shock height and speed
using a coronal electron-density model.

\section{The analysis chain}

The analysis uses four tools, normally applied in this order:
\begin{steps}
  \item \textbf{Burst isolation} restricts the data to the emission of interest
    (Section~\ref{sec:isolation}).
  \item \textbf{Maximum intensities} extracts the frequency of peak emission in each time
    channel, or tracks the burst ridge (Sections~\ref{sec:max-intensities}
    and~\ref{sec:ridge-tracking}).
  \item \textbf{Outlier removal} discards points that do not belong to the burst
    (Section~\ref{sec:outliers}).
  \item \textbf{The Burst Analyzer} fits the backbone and derives the drift rate, shock
    speed and shock height (Section~\ref{sec:analyzer}).
\end{steps}
Each tool works on the spectrum as currently displayed: background-subtracted,
RFI-cleaned or isolated. Figure~\ref{fig:type-ii-spectrum} shows a background-subtracted
spectrum containing a Type~II burst.

\screenshot{analysis/loaded_fits}{A background-subtracted dynamic spectrum containing a Type~II burst.}{fig:type-ii-spectrum}{Main window showing a background-subtracted spectrum with a Type II burst.}

\section{Burst isolation}
\label{sec:isolation}
\index{burst isolation}\index{lasso}

The \gui{Isolate Burst}\index{burst isolation!Isolate Burst tool} toolbar tool (tool~9) draws a freehand lasso mask around the
emission region of a burst. Background subtraction must be applied first.
\begin{steps}
  \item Click \gui{Isolate Burst}.
  \item Press the mouse button, drag a loop around the burst, and release.
\end{steps}
Only the data inside the loop are retained; everything outside is set to zero, which is the
background level of the background-subtracted spectrum. This removes unrelated emission and
residual noise from the analysis (Figure~\ref{fig:isolated}). Isolation can be undone\index{undo} with
\kbd{Ctrl+Z}. The plot
title changes to \gui{-Isolated Burst}. The mask follows the drawn path on the displayed
spectrum. \gui{Reset Selection}\index{Reset Selection} (tool~14) clears the selection.

\screenshot{analysis/isolated_burst}{An isolated Type~II burst lane.}{fig:isolated}{Main window after burst isolation.}

\section{The Maximum Intensities window}
\label{sec:max-intensities}
\index{maximum intensities}

The Maximum Intensities window (Figure~\ref{fig:max-intensities}) shows, for each time
channel of the displayed spectrum, the frequency at which the intensity is highest. These
points trace the burst backbone that the Analyzer fits. Open it with the
\gui{Plot Maximum Intensities}\index{maximum intensities!toolbar tool} toolbar tool (tool~10) or
\menu{Analysis > Maximum Intensities > Open Maximum Intensities}. When the maxima are taken
from an isolated burst, time channels with no emission inside the mask are excluded.

\screenshot{analysis/max_intensities}{The Maximum Intensities window. Without isolation, a persistent RFI channel near 89~MHz produces spurious points.}{fig:max-intensities}{Maximum Intensities window.}

\begin{table}[!htbp]
  \centering
  \caption{Controls of the Maximum Intensities window.}
  \label{tab:max-controls}
  \small
  \begin{tabularx}{\linewidth}{@{}P{0.3\linewidth}L@{}}
    \toprule
    \thead{Control} & \thead{Function} \\
    \midrule
    \gui{Select Outliers}\index{outlier removal!manual} & Draw a lasso around points to remove. \\
    \gui{Remove Outliers} & Remove the selected points. \\
    \gui{Track Ridge}\index{ridge tracking!Track Ridge} & Replace the per-column maxima by an automatically tracked ridge (Section~\ref{sec:ridge-tracking}). \\
    \gui{Pick Start}\index{ridge tracking!Pick Start} & Click a point on the burst from which ridge tracking starts. \\
    \gui{Fundamental} / \gui{Harmonic} & Whether the lane is fundamental or harmonic plasma emission (Section~\ref{sec:harmonic}). \\
    \gui{Analyze Burst}\index{Burst Analyzer!opening} & Open the Burst Analyzer with the current points. \\
    \menu{File > Save As} & Save the points as CSV: time channel, time (s) and frequency (MHz). \\
    \menu{File > Export As} & Save the plot as a figure (PNG, PDF, EPS, SVG, TIFF or JPG). \\
    \menu{Edit > Reset All} & Clear the selection, the ridge start point and the stored Analyzer state. \\
    \menu{Edit > Restore Per-Column Maxima}\index{ridge tracking!restore maxima} & Bring back the points replaced by ridge tracking. \\
    \menu{Analyze > Open Analyzer} & Open the Burst Analyzer. \\
    \menu{Analyze > Ridge Tracking Settings...}\index{ridge tracking!settings} & Parameters of the ridge tracker (Table~\ref{tab:ridge-settings}). \\
    \bottomrule
  \end{tabularx}
\end{table}

The status bar at the bottom of the window reports the result of each action, for example
the number of selected points.

\section{Ridge tracking}
\label{sec:ridge-tracking}
\index{ridge tracking}

Taking the brightest channel in every time column picks up RFI, a second burst or noise,
leaving scattered points to remove by hand. \gui{Track Ridge} follows the burst itself
instead:
\begin{itemize}
  \item Tracking starts from the brightest point of the burst, or from a point chosen with
    \gui{Pick Start} (marked with a star), and continues forwards and backwards in time,
    searching only a few channels around the ridge position in the previous column.
  \item A channel that is bright for most of the file is treated as RFI and is not chosen
    as a starting point. An error is reported if a picked start point is not above the
    background.
  \item Short dropouts in the burst are bridged. After a gap, tracking continues only once
    two consecutive columns show signal, so a single noise spike cannot extend the ridge.
    Tracking stops once the signal has stayed below the threshold for longer than the
    allowed gap.
  \item Each peak is refined to sub-channel frequency.
\end{itemize}
The tracked points replace the per-column maxima; outlier removal, the
fundamental/harmonic choice and \gui{Analyze Burst} work on them as before.
\menu{Edit > Restore Per-Column Maxima} brings the original points back.

\begin{table}[!htbp]
  \centering
  \caption{Ridge Tracking Settings.}
  \label{tab:ridge-settings}
  \small
  \begin{tabularx}{\linewidth}{@{}P{0.3\linewidth}LP{0.14\linewidth}@{}}
    \toprule
    \thead{Setting} & \thead{Description} & \thead{Default} \\
    \midrule
    Search window (channels) & Maximum movement of the ridge between two consecutive time columns (1--100). & 3 \\
    Threshold (\(\sigma\)) & Minimum height of a peak above the background, in robust standard deviations (0--50). & 3.0 \\
    Allowed gap (columns) & Number of consecutive columns below the threshold before tracking stops (0--1000). & 8 \\
    Only follow drift to lower frequency & Restricts the ridge to move toward lower frequencies with time, as in Type~II and Type~III bursts. & Off \\
    \bottomrule
  \end{tabularx}
\end{table}

\screenshot[0.5\linewidth]{dialogs/ridge_tracking_settings}{The \gui{Ridge Tracking Settings} dialog.}{fig:ridge-settings}{The Ridge Tracking Settings dialog with its four fields and OK/Cancel.}

\section{Outlier removal}
\label{sec:outliers}
\index{outlier removal}

Outliers are removed manually in the Maximum Intensities window:
\begin{steps}
  \item Click \gui{Select Outliers} and draw a lasso around the unwanted points
    (Figure~\ref{fig:outliers}).
  \item Click \gui{Remove Outliers}.
\end{steps}

\screenshot{analysis/outlier_removal}{Selecting outliers with the lasso in the Maximum Intensities window.}{fig:outliers}{Outlier selection in the Maximum Intensities window.}

When \menu{Processing > Maximum Intensity > Auto-Clean Isolated Burst Outliers}\index{outlier removal!automatic} is ticked
(the default), outliers are also removed automatically whenever the maxima are taken from
an isolated burst. The manual tools remain available in either case.

\section{Fundamental and harmonic emission}
\label{sec:harmonic}
\index{harmonic emission}

Type~II emission occurs at the local plasma frequency (fundamental) and at twice that
frequency (harmonic). Select \gui{Harmonic} when the fitted lane is the harmonic. The
observed frequencies and drift rates are then divided by two before the shock parameters
are calculated, so that heights and speeds always refer to the plasma frequency. Saved
summaries keep both the converted values and the observed reference values.

\section{The Burst Analyzer}
\label{sec:analyzer}
\index{Burst Analyzer}

The Burst Analyzer (Figure~\ref{fig:analyzer}) fits the burst backbone with a power law and
derives the frequency drift rate, shock speed and shock height, with the goodness of fit.
Open it with \gui{Analyze Burst} or \menu{Analyze > Open Analyzer} in the Maximum
Intensities window.

\screenshot{analysis/analyzer}{The Burst Analyzer window with the best fit, shock parameters and density-model comparison.}{fig:analyzer}{Burst Analyzer window.}

\begin{table}[!htbp]
  \centering
  \caption{Controls and displays of the Burst Analyzer.}
  \label{tab:analyzer}
  \small
  \begin{tabularx}{\linewidth}{@{}P{0.29\linewidth}L@{}}
    \toprule
    \thead{Control or display} & \thead{Function} \\
    \midrule
    \gui{Maximum Intensities} & Plot the extracted points. \\
    \gui{Best Fit} & Fit the points with a power law and plot the fit. \\
    \gui{t₀} & Time origin of the power law (Section~\ref{sec:t0}). \\
    \gui{Density model} & Coronal density model for heights and speeds (Section~\ref{sec:density-models}). \\
    \gui{Fold-number}\index{fold number} & Multiplier 1--4 applied to the model density; press \gui{Calculate} after changing it. \\
    \gui{Best Fit Equation} & The fitted power law. \\
    \gui{Fit Metrics} & \(R^2\) and RMSE of the fit. \\
    \gui{Shock Parameters} & Derived quantities with uncertainties (Table~\ref{tab:shock-params}). \\
    \gui{Density model comparison}\index{density model!comparison table} & Initial and average shock speed and height from all five models (Section~\ref{sec:model-comparison}). \\
    \gui{Save Graph}\index{Burst Analyzer!Save Graph} & Save the plot on show as a figure. \\
    \gui{Save Data}\index{Burst Analyzer!Excel export} & Save the results to an Excel workbook (Section~\ref{sec:analyzer-export}). \\
    \gui{Existing Excel File} & Append to an existing workbook instead of creating a new one. \\
    \gui{Extra Plots}\index{Burst Analyzer!extra plots} and \gui{Plot} & Additional plots (Section~\ref{sec:extra-plots}). \\
    \bottomrule
  \end{tabularx}
\end{table}

\subsection{Power-law fit and the time origin}
\label{sec:t0}
\index{power law}\index{time origin}

The backbone is fitted as
\begin{equation}
  f(t) = a\,(t - t_0)^{-b},
  \label{eq:power-law}
\end{equation}
where \(f\) is the frequency in MHz, \(t\) the time in seconds, and \(a\) and \(b\) the fitted
parameters. The drift rate at each point follows from the derivative,
\(\mathrm{d}f/\mathrm{d}t = -ab\,(t-t_0)^{-b-1}\). On the \gui{Best Fit} graph the points are
drawn as filled black squares and the fitted power law as a red curve, with a boxed legend.

Because a power law depends on where \(t=0\) lies, the fitted exponent, drift rate and shock
parameters depend on the chosen time origin. \(t_0\) is selected from the \gui{t₀} list next
to \gui{Best Fit}:
\begin{description}
  \item[File start] measures time from the start of the loaded data (the default).
  \item[Burst onset] places \(t_0\) one time sample before the first fitted point, so that
    every point lies on the curve.
  \item[Custom time] uses a time that you enter, in UT when the start time of the file is
    known and otherwise in seconds from the start of the file.
\end{description}
Changing \(t_0\) re-runs the fit automatically. The graph keeps the data's own time axis,
and the equation shows \((x - t_0)\) whenever \(t_0\) is not the file start.

\subsection{Shock parameters}

Table~\ref{tab:shock-params} defines the quantities in the \gui{Shock Parameters} block. The
heading names the density model and fold in use, for example
\gui{Shock Parameters (Newkirk 1-fold)}\index{density model!Newkirk}.

\begin{table}[!htbp]
  \centering
  \caption{Shock parameters reported by the Burst Analyzer.}
  \label{tab:shock-params}
  \small
  \begin{tabularx}{\linewidth}{@{}P{0.27\linewidth}L@{}}
    \toprule
    \thead{Quantity} & \thead{Definition} \\
    \midrule
    Average Frequency & Mean frequency of the maximum-intensity points (MHz), with its standard error. \\
    Average Drift Rate & Mean of the fitted drift rate along the backbone (MHz/s). \\
    Starting Frequency\index{shock parameters!starting frequency} & The 90th percentile of the fitted frequencies, taken as the start of the burst (MHz). \\
    Initial Shock Speed\index{shock parameters!initial speed} & Shock speed at the starting frequency (km/s). \\
    Initial Shock Height & Heliocentric height at the starting frequency (\(R_\odot\)). \\
    Average Shock Speed\index{shock parameters!average speed} & Mean shock speed along the backbone (km/s), with its standard error. \\
    Average Shock Height & Mean height along the backbone (\(R_\odot\)), with its standard error. \\
    \bottomrule
  \end{tabularx}
\end{table}

\begin{background}[Background: from drift rate to shock speed]
  The emission frequency is identified with the local plasma frequency,
  \(f_p \approx 8.98\times10^{-3}\sqrt{n_e}\)~MHz for an electron density \(n_e\) in
  cm\(^{-3}\). A density model \(n_e(r)\) therefore maps each frequency to a heliocentric
  height \(r\), and the drift rate to a radial speed,
  \begin{equation*}
    v = \left|\frac{\mathrm{d}f}{\mathrm{d}t}\right| \left|\frac{\mathrm{d}r}{\mathrm{d}f}\right|,
    \qquad \frac{\mathrm{d}r}{\mathrm{d}f} = \frac{2\,n_e}{f\,|\mathrm{d}n_e/\mathrm{d}r|}.
  \end{equation*}
  Uncertainties are propagated linearly from the drift-rate and frequency errors. The
  results are only as good as the density model, which describes an average corona and
  can differ from the actual medium by large factors; the fold number scales the model
  density to account for denser structures such as streamers.
\end{background}

\subsection{Density models}
\label{sec:density-models}
\index{density model}

Shock heights and speeds are calculated from one of five coronal electron-density models
(Table~\ref{tab:density-models}). Changing the model recalculates the shock parameters
immediately. The fold number (1--4) multiplies the model density; press \gui{Calculate}
after changing it. The formulas are given in Appendix~\ref{app:formulas}.

\begin{table}[!htbp]
  \centering
  \caption{Coronal electron-density models.}
  \label{tab:density-models}
  \small
  \begin{tabularx}{\linewidth}{@{}P{0.25\linewidth}L@{}}
    \toprule
    \thead{Model} & \thead{Reference} \\
    \midrule
    Newkirk (default) & \textcite{newkirk1961} \\
    Saito\index{density model!Saito} & \textcite{saito1977} \\
    Leblanc\index{density model!Leblanc} & \textcite{leblanc1998} \\
    Baumbach\index{density model!Baumbach--Allen}--Allen & \textcite{baumbach1937}; \textcite{allen1947} \\
    Mann & \textcite{mann1999}\index{density model!Mann} \\
    \bottomrule
  \end{tabularx}
\end{table}

Models other than Newkirk are evaluated between the photosphere (\(1\,R_\odot\)) and 1~AU.
A frequency that a model cannot reach in this range --- for example above about 116~MHz
for the 1-fold Saito model --- has no height, and is shown as a dash rather than an
extrapolated value.

\subsection{Density-model comparison}
\label{sec:model-comparison}

The \gui{Density model comparison} table below the shock parameters lists the initial shock
speed \(v_0\), initial height \(h_0\), average speed \(\bar v\) and average height \(\bar h\)
obtained with every model at the current fold, with the selected model in bold. It shows at
a glance how strongly the derived kinematics depend on the assumed corona.

\subsection{Additional plots}
\label{sec:extra-plots}

Choose an entry in \gui{Extra Plots} and click \gui{Plot} to show:
\gui{Shock Speed vs Shock Height}, \gui{Shock Speed vs Frequency} or
\gui{Shock Height vs Frequency}. These plots use the selected density model.

\subsection{Saving graphs and data}
\label{sec:analyzer-export}
\index{Excel export}

\gui{Save Graph} writes whichever plot is on show --- maximum intensities, best fit or an
additional plot --- as an OriginPro-style graph on a white page in PNG, PDF, EPS, SVG, TIFF
or JPG format.

\gui{Save Data} writes one row of results to an Excel workbook (\file{.xlsx}). With
\gui{Existing Excel File} ticked, you choose a workbook and the row is appended, which builds
an event table over many analyses. The columns are listed in Table~\ref{tab:excel-columns}.

\begin{table}[!htbp]
  \centering
  \caption{Columns of the Burst Analyzer Excel export.}
  \label{tab:excel-columns}
  \small
  \begin{tabularx}{\linewidth}{@{}rL@{}}
    \toprule
    \thead{No.} & \thead{Column} \\
    \midrule
    1--2 & Observation date; station \\
    3--5 & Best-fit equation; \(R^2\); RMSE \\
    6--9 & Average frequency and error; average drift rate and error \\
    10--11 & Starting frequency and error \\
    12--15 & Initial shock speed and error; initial shock height and error \\
    16--19 & Average shock speed and error; average shock height and error \\
    20 & Absolute value of the average drift rate \\
    21 & Density model \\
    22 & \(t_0\) in seconds from the start of the file \\
    \bottomrule
  \end{tabularx}
\end{table}

The density model and \(t_0\) are also saved in projects and named in project reports.
Projects saved before these options existed open with the Newkirk model and the file start
as \(t_0\), matching how they were calculated.
